\section{Course Schedule}
\label{sec:cses-course-schedule}

\problemstatement{Given $n$ courses and $m$ prerequisite relations (a
directed edge $a\to b$ means $a$ must be taken before $b$), output a valid
order to complete all courses, or \texttt{IMPOSSIBLE} if the prerequisites
are cyclic.}

\begin{patternbox}
Ordering tasks under ``before'' constraints is \textbf{topological sort}.
Kahn's algorithm repeatedly outputs a node with no remaining prerequisites
(in-degree $0$); if some nodes never reach in-degree $0$, a cycle exists and
no order is possible.
\end{patternbox}

\subsection*{Kahn's in-degree queue}

Compute each node's in-degree (number of prerequisites). Put all in-degree-
$0$ nodes in a queue. Repeatedly pop one, append it to the order, and
decrement the in-degree of its successors; any successor reaching $0$ joins
the queue. If the produced order contains all $n$ nodes, it is a valid
topological order; if fewer, the leftover nodes form a cycle and the answer
is \texttt{IMPOSSIBLE}.

\begin{ideabox}[In-Degree Zero Means ``Ready Now'']
A node can be scheduled once all its prerequisites are done -- i.e.\ its
in-degree has dropped to $0$. Draining nodes in this readiness order both
produces a valid schedule and detects a cycle (some nodes never become
ready).
\end{ideabox}

\subsection*{Algorithm}

\begin{enumerate}
  \item Build the graph and in-degrees; queue all in-degree-$0$ nodes.
  \item Pop a node, output it, decrement successors' in-degrees, enqueue
        new zeros.
  \item If all $n$ nodes were output, print the order; else
        \texttt{IMPOSSIBLE}.
\end{enumerate}

\subsection*{Dry run}

\dryrun{Edges $1\!\to\!2$, $1\!\to\!3$, $3\!\to\!2$.}

\begin{itemize}
  \item In-degrees: $1{:}0,2{:}2,3{:}1$. Start with $1$; then $3$ becomes
        $0$; then $2$.
  \item Order $1,3,2$.
\end{itemize}

\subsection*{C++ implementation}

\begin{lstlisting}
#include <bits/stdc++.h>
using namespace std;

int main() {
    int n, m;
    cin >> n >> m;
    vector<vector<int>> adj(n + 1);
    vector<int> indeg(n + 1, 0);
    for (int i = 0; i < m; i++) {
        int a, b; cin >> a >> b;
        adj[a].push_back(b);
        indeg[b]++;
    }

    queue<int> q;
    for (int i = 1; i <= n; i++) if (indeg[i] == 0) q.push(i);

    vector<int> order;
    while (!q.empty()) {
        int u = q.front(); q.pop();
        order.push_back(u);
        for (int v : adj[u]) if (--indeg[v] == 0) q.push(v);
    }

    if ((int)order.size() != n) { cout << "IMPOSSIBLE\n"; return 0; }
    for (int v : order) cout << v << " ";
    cout << "\n";
    return 0;
}
\end{lstlisting}

\begin{complexitybox}
\textbf{Time Complexity.} $O(n+m)$. \textbf{Space Complexity.} $O(n+m)$.
\end{complexitybox}

\begin{takeawaybox}
Scheduling under precedence is topological sort; Kahn's in-degree queue both
produces the order and detects cyclic (hence infeasible) constraints when
fewer than $n$ nodes emerge. This is the standard dependency-resolution
algorithm (build systems, task runners).
\end{takeawaybox}
